\section{Auswertung}
\label{sec:Auswertung}

% Fehlerrechnung
Die in diesem Kapitel durchgeführten Fehlerrechnungen wurden mittels der Gaußschen Fehlerfortpflanzung 
\begin{equation}
    \sigma_f = \sqrt{\left( \frac{\partial f}{\partial x_1} \sigma_{x_1} \right)^2 + \left( \frac{\partial f}{\partial x_2} \sigma_{x_2} \right)^2 + ...}
\end{equation}
durchgeführt, sowie den Formeln 
\begin{equation}
    \overline{x} = \frac{1}{n} \sum^n_{i=1}x_i . 
    \label{eq:mittelwert}
\end{equation}
für den Mittelwert und 
\begin{equation}
    \sigma_{\overline{x}} = \frac{\sigma_x}{\sqrt{n}} = \sqrt{\frac{1}{n(n-1)}  \sum^n_{i=1} \left( x_i - \overline{x}\right)^2 } .
    \label{eq:abweichung}
\end{equation}
für die Standardabweichung einer Messreihe $\left( x_1, ... , x_n  \right)$.
Der Einfachheit halber wird hierbei das Python-Paket \textit{uncertainties} verwendet. \cite{uncertainties} \\
\hspace{1cm}
% Ende Fehlerrechnung



\clearpage